% Use utf-8 encoding for foreign characters
\usepackage[utf8]{inputenc}

% Setup for fullpage use
% \usepackage{fullpage}

% Running Headers and footers
% \usepackage{fancyhdr}
% \usepackage{newhead}

% More symbols
\usepackage{amsmath, amssymb, amsthm, amstext, latexsym}

% Surround parts of graphics with box
% \usepackage{boxedminipage}
% \usepackage{psboxit}
% \PScommands

\usepackage{framed}

\usepackage[a4paper]{geometry}            % See geometry.pdf to learn the layout options. There are lots.
\usepackage[parfill]{parskip}    % Activate to begin paragraphs with an empty line rather than an indent

% This is now the recommended way for checking for PDFLaTeX:
\usepackage{ifpdf}

\ifpdf
  \usepackage[pdftex]{graphicx}
  \pdfcompresslevel=9
  \DeclareGraphicsExtensions{.pdf,.png,.jpg} 
  \DeclareGraphicsRule{.tif}{png}{.png}{`convert #1 `dirname #1`/`basename #1 .tif`.png}
\else
  \usepackage{graphicx}
\fi

\usepackage{epstopdf}

\usepackage[usenames,dvipsnames]{xcolor}

\usepackage[font=small,labelfont=bf,position=bottom]{caption}
% \usepackage{makeidx}
\usepackage{multicol}
\usepackage{multirow}
\usepackage[sectionbib,square]{natbib}

\usepackage{wrapfig} % Allow text to wrap figures
\usepackage{subfig} % Allow subfigures in a figure

\usepackage{flushend}
\usepackage{rotating}
\usepackage{url}
\usepackage{varioref}
% 
% \usepackage{exercise}

% ======================
% = Playing with Fonts =
% ======================

\usepackage[OT1]{fontenc}
\usepackage{babel}
\usepackage{array}
\usepackage{amsmath}
\usepackage{slantsc}
\usepackage{yfonts}

\usepackage{times}
\usepackage{tgbonum}  %% Use the extended bookman font
\usepackage{aurical}
\renewcommand*\ttdefault{txtt}

\usepackage{morse}

% \usepackage{mybox}
% \definecolor{shadecolor}{rgb}{0.9,0.9,0.9}

%\usepackage{noteslang}
%\usejava
%\usecsharp

% =========================
% = Packages for graphics =
% =========================
\usepackage{tikz}
\usepackage{pgf}

\input{preamble/tikz-style}

% =======================
% = Drop Caps and Fonts =
% =======================

\usepackage{lettrine}

% ============
% = Glossary =
% ============
\usepackage[xindy]{glossaries}

% ==============================================
% = Packages required for Source Code Listings =
% ==============================================

\usepackage{listings}
\usepackage{upquote}
\usepackage{fancyvrb}
\usepackage{pdflscape}

\lstdefinelanguage[FPC]{Pascal}[Standard]{Pascal}
  {morecomment=[l]//,
   morekeywords={uses,external,cdecl,export,forward,global,module,nil,operator,overload,%
      priority,sum,type,use,dispose,mark,page,release,%
      eof,eoln,ord,pos,pred,rval,succ,result,interface,implementation,unit}%
  }%

\lstset{keywordstyle=\color{Blue}}
\lstset{commentstyle=\color{OliveGreen}}
\lstset{stringstyle=\color{BrickRed}}
\lstset{basicstyle=\ttfamily\small}
\lstset{captionpos=t}
% \lstset{frame=trbl}
\lstset{fancyvrb=true}
\lstset{numberstyle=\footnotesize}
\lstset{showstringspaces=false}
\lstset{tabsize=2}
\lstset{captionpos=b}

\newcommand{\straightcode}[1]
{
  \begin{tikzpicture}
    \draw[very thick,color=gray] (0,-0.15) -- (0,0) -- (\textwidth,0) -- (\textwidth,-0.15);
  \end{tikzpicture}
  
  {#1}
  
  \begin{tikzpicture}
    \draw[very thick,color=gray] (0,0.15) -- (0,0) -- (\textwidth,0) -- (\textwidth,0.15);
  \end{tikzpicture}
}


\newcommand{\pascode}[3]
{
  \lstset{language=[FPC]Pascal}
  \lstinputlisting[caption={#2},label={#1}]{#3}
}

\newcommand{\mpascode}[4]
{
  \lstset{language=[FPC]Pascal}
  \lstinputlisting[caption={#2},label={#1},multicols=#4]{#3}
}


\newcommand{\passnipet}[1]
{
  \lstset{language=[FPC]Pascal}
  \lstinline&#1&
}


\newcommand{\ccode}[3]
{
  \lstset{language=C}
  \lstinputlisting[caption={#2},label={#1}]{#3}
}

\newcommand{\mccode}[4]
{
  \lstset{language=C}
  \lstinputlisting[caption={#2},label={#1},multicols=#4]{#3}
}


\newcommand{\csnipet}[1]
{
  \lstset{language=C}
  \lstinline^#1^
}

\newcommand{\bashcode}[3]
{
  \lstset{language=bash}
  \lstinputlisting[caption={#2},label={#1}]{#3}
}

\newcommand{\bashsnipet}[1]
{
  \lstset{language=bash}
  \lstinline^#1^
}


\captionsetup[lstlisting]{labelfont=bf} 

\addtolength{\evensidemargin}{-0.50in}
\addtolength{\textwidth}{0.50in}

\addtolength{\topmargin}{-0.25in}
\addtolength{\textheight}{1.25in}

\usepackage{phdthesis}

% ==================
% = MyBox commands =
% ==================

\usepackage{lipsum,framed,color}

\definecolor{shadecolor}{rgb}{1,0.8,0.3}
\newlength{\boxwidth}
\newsavebox{\boxcontainer}

\newcommand\mybox[1]{%
  \begin{center}%
    \fcolorbox{black}{shadecolor}{%
      \begin{minipage}[t]{\dimexpr\textwidth-2\fboxsep-2\fboxrule}%
        #1%
      \end{minipage}}%
  \end{center}}

\newenvironment{Mybox}%
  {\begin{lrbox}{\boxcontainer}%
     \begin{minipage}{\dimexpr\textwidth-2\fboxsep-2\fboxrule}}%
  {\end{minipage}\end{lrbox}%
   \begin{center}%
     \fcolorbox{black}{shadecolor}{\usebox{\boxcontainer}}
   \end{center}}

\newenvironment{MyBox}{%
  \def\FrameCommand{\fcolorbox{black}{shadecolor}}%
  \MakeFramed{\setlength{\hsize}{\dimexpr\textwidth-2\fboxsep-2\fboxrule} \FrameRestore}}%
{\endMakeFramed}

% Settings for \mybox and Mybox
\setlength{\fboxrule}{0pt}
\setlength{\fboxsep}{9pt}

% Additional settings for MyBox
\setlength{\FrameRule}{\fboxrule}
\setlength{\FrameSep}{\fboxsep}

% ===========================
% = Bookmarks for PDF Index =
% ===========================
\usepackage[bookmarks=true,colorlinks=true,linkcolor=blue,plainpages=false,pdfpagelabels]{hyperref}

\input{preamble/commands}

% If you want to generate a toc for each chapter (use with book)
\usepackage{minitoc}

\setcounter{totalnumber}{50}
\setcounter{topnumber}{50}
\setcounter{bottomnumber}{50}

\renewcommand\floatpagefraction{.9}
\renewcommand\topfraction{.9}
\renewcommand\bottomfraction{.9}
\renewcommand\textfraction{.1}

% ============
% = Metadata =
% ============

\title{The Programming Arcana}
\author{Andrew Cain}
\date{\today}
